\section{Heaps and Priority Queues}

Priority queues support the operations:
\begin{enumerate}
    \setlength{\itemsep}{0em}
    \item \textbf{Max()}: Return elment with largest key.
    \item \textbf{ExtractMax()}: Remove and return element with largest key.
    \item \textbf{IncreaseKey(x, k)}: Increase the key of element $x$ to $k$.
    \item \textbf{Insert(x)}: Insert element $x$ into the priority queue.
\end{enumerate}

\subsection{Heaps}
Heaps are \emph{almost complete} binary trees\footnote{All nodes have two children except for the last level, which is filled from left to right}
where each node has a key that is greater than or equal to the keys of its children.

\subsubsection{Operations}
\textbf{BubbleUp(x)}: If the key of $x$ is greater than the key of its parent, swap $x$ with its parent and repeat until the heap property is restored.

\textbf{BubbleDown(x)}: If the key of $x$ is less than the key of one of its children, swap $x$ with the child with the largest key and repeat until the heap property is restored.

\timeComplexity{\log n}{1}

\subsubsection{Representations}

\textbf{Linked list}:
Each node has pointers to its parent and children.
Allows arbitrary size. 

\textbf{Array}:
The root is at index 1, the left child of $i$ is $2i$, the right child $2i + 1$, and the parent is $\lfloor i / 2 \rfloor$.
More space efficient than linked list.

\subsection{Building heaps}
To turn an array into a heap we use \emph{Bottom-up construction}.
For the first $[1, \lfloor n / 2 \rfloor]$ elements, we call \textbf{BubbleDown(i)}, starting from $\lfloor n / 2 \rfloor$ and going down to 1.

\timeComplexity{n}{1}

\subsection{Using heaps for priority queues}
Below is baed on heaps as arrays. Logic is similar for linked lists.
\begin{python}
def max(heap):
    return heap[1]

def extract_max(heap):
    max_element = heap[1]
    heap[1] = heap[length] # Move last element to root.
    length -= 1 # Remove last element.
    bubble_down(heap, 1) # Restore heap property.
    return max_element

def increase_key(heap, i, k):
    heap[i] = k
    bubble_up(heap, i) # Restore heap property.

def insert(heap, k):
    length += 1
    heap[length] = k
    bubble_up(heap, length) # Restore heap property.
\end{python}

\subsection{HeapSort}
To sort an array, build a heap and repeatedly run \verb|extractMax()|.

This has the advantage of being an in-place sorting algorithm.

\timeComplexity{n \lg n}